% Discussion update: mechanism interpretation
% Paste into Discussion chapter, or as a subsection after E1/E2 results if a separate discussion chapter is not yet used.

\section{Mechanistic interpretation of the early reorganisation}
\label{sec:mechanism-interpretation}

The E1 results do not by themselves identify a unique mechanism. They do, however, constrain what a plausible mechanism must explain. Across model sizes, several spectral signatures concentrate in the same early part of training: leading-subspace instability, MP-outlier emergence in rectangular matrices, stable-rank collapse, and heavy-tail changes. These signatures suggest a transition from an initially near-random regime, where weight matrices are still close to their random initial spectra, to a regime in which a small number of data-aligned directions begin to dominate.

A useful interpretation is therefore a two-part ``feature-subspace formation and consolidation'' mechanism. In the first part, training accumulates signal from high-frequency, high-signal-to-noise regularities in the corpus. These regularities break the approximate random-matrix structure of the initial weights, producing spectral outliers and rapid rotation of leading singular subspaces. In the second part, the model compresses much of the early learning into a lower-dimensional set of dominant directions. This appears as stable-rank and effective-rank collapse. The fact that MP-outlier and subspace-stability indicators often move earlier than rank-compression indicators is consistent with this sequence: dominant directions first separate from the random bulk, and broader rank compression follows as those directions are amplified and stabilised.

This interpretation also explains why the phase boundary is module-staggered. Attention query/key/value matrices are directly involved in routing and contextual selection, so their leading subspaces can begin reorganising early as token co-occurrence and positional regularities become useful. MLP projections, especially the expansion projection, may reorganise later because they encode and transform features once the routing structure has begun to stabilise. The empirical result is therefore not a single global switch, but a coordinated sequence of component-wise transitions.

The interpretation remains deliberately conservative. The E1 metrics are weight-space diagnostics, not direct circuit reconstructions. The claim is not that the model has learned a specific linguistic algorithm at step 2000. The claim is that the weights undergo a reproducible transition from random-like high-rank structure to data-aligned lower-rank structure, and that this transition provides a plausible substrate for stage-dependent plasticity. E3 tests the consequential part of this interpretation: whether signals introduced during or after this window are retained differently.

This mechanism also explains why a universal-looking window can appear across model sizes. Pythia models share architecture family, training data, data order, and checkpoint schedule, so early high-signal corpus regularities are presented at comparable token counts across scales. Larger models need not have identical transition sharpness, but if the same high-SNR statistics drive early subspace formation, the broad ordering of phases can be preserved. This is the sense in which the observed dynamics can be universal: not because every model has the same exact boundary, but because models trained under the same data and optimisation regime pass through a similar sequence of spectral events.
